% =============================================================
% STRUCTURED ABSTRACT for JAMIA Research and Applications
% (≤250 words; five headings:
%  Objective / Materials and Methods / Results / Discussion /
%  Conclusion).
% This file expands inside \begin{abstract}...\end{abstract}.
% =============================================================

\noindent\textbf{Objective.}
In cancer knowledge-base (KB) curation, models are typically
selected by public benchmark scores. We test whether benchmark F1
predicts KB-surfacing yield in a pre-registered factorial study.

\noindent\textbf{Materials and Methods.}
We trained 310 models across two waves: a schema-comparison wave
(4 encoders \(\times\) 3 schemas \(\times\) 10 seeds = 120 runs)
selecting the operational relation schema, and a configuration
wave (3 encoders \(\times\) 3 training schedules \(\times\) 20
seeds + a 10-seed RoBERTa-base reference = 190 runs) under that
schema. Each model was evaluated against 162 curator-validated
CIViC oncology assertions. Analyses comprised
pre-registered paired-difference contrasts on three KB-audit
metrics, a variance decomposition by design lever, and the
rank-inversion rate among benchmark-tied configuration pairs.

\noindent\textbf{Results.}
An entity-pair schema improved KB surfacing over a coarser family
schema (\(\Delta=+0.117\), 95\% paired-bootstrap CI [+0.045, +0.194]).
Multi-corpus training improved out-of-domain generalisation on
BC5CDR drug--disease F1 (\(\Delta=+0.139\), Cohen's \(d=2.04\)).
Encoder, schedule, and their interaction explained about five
times more variance in KB argmax accuracy than in BioRED ex-NEG F1
(\(R_B=0.21\); cluster-bootstrap 95\% CI [0.028, 0.990]), with the
training-schedule term contributing most of the gap.
Among 18 benchmark-tied configuration pairs the rank-inversion rate
was 0.50 (95\% CI [0.14, 0.83]).

\noindent\textbf{Discussion.}
Benchmark F1 and KB yield diverged most in the training-schedule
arm: half of benchmark-tied configuration pairs reversed their KB
ranking. Wide confidence intervals at the realised cell counts
limit directional claims for the coupling-slope family.

\noindent\textbf{Conclusion.}
Benchmark scores alone are insufficient for oncology KB curation
pipelines. We recommend pairing benchmark scoring with a
lightweight KB-anchor audit on 100--200 curator-validated
assertions; the audit code and pre-registration are released.
